\section{Crisis Director Letter}

\begin{chairletter}[IEUMUN 2026 Crisis Director]{Noelia Álvarez Iglesias}
Dear Delegates,

It is my pleasure to welcome you to the Crisis Committee of IEUMUN 2026. Whether this is your first experience with crisis or you have already spent countless sessions writing directives, debating strategy, and navigating unexpected updates, I hope this committee offers you a challenging, immersive, and memorable experience.

This year you will enter one of the Cold War’s most intricate and controversial theatres: Latin America in 1975. Rather than focusing on a single government or institution, this committee has been designed as a quadcameral crisis, bringing together four distinct perspectives whose interests rarely align but whose decisions are deeply interconnected. Throughout the conference, governments, intelligence agencies, and revolutionary movements will operate simultaneously, shaping the same reality from different angles. Every directive, negotiation, and strategic calculation has the potential to alter the balance of power, often in ways that extend beyond your own cabinet.

Unlike traditional committees, crisis demands adaptability above all else. Success will not simply come from having the strongest policy or the loudest speeches, but from anticipating consequences, working effectively with your cabinet, and responding to developments that cannot always be predicted. While history provides the starting point, the future of this committee will ultimately be determined by your choices.

Above all, I encourage you to embrace the experience. Be creative, think critically, challenge assumptions, and do not be afraid to take risks. Some of the most memorable moments in crisis emerge from ideas that no one expected, and there is rarely a single correct solution to the problems you will face.

On a more personal note, crisis committees hold a special place for me. Last year, IEUMUN was my first experience as a backroomer, and it quickly became one of the most rewarding experiences I have ever had in MUN. Watching delegates navigate uncertainty, devise ambitious strategies, and completely reshape the course of history showed me what makes crisis such a unique format. It was that experience that inspired me to return this year as Crisis Director. Together with my team, we have worked to create a committee that will challenge and excite you, and I sincerely hope it offers you an unforgettable experience. I encourage you to embrace every opportunity, think boldly, and enjoy every moment.

I look forward to seeing the story that you will create.

\signoff[Best regards,]{Noelia Álvarez Iglesias}
\end{chairletter}
